% Options for packages loaded elsewhere
\PassOptionsToPackage{unicode}{hyperref}
\PassOptionsToPackage{hyphens}{url}
%
\documentclass[
]{article}
\usepackage{amsmath,amssymb}
\usepackage{iftex}
\ifPDFTeX
  \usepackage[T1]{fontenc}
  \usepackage[utf8]{inputenc}
  \usepackage{textcomp} % provide euro and other symbols
\else % if luatex or xetex
  \usepackage{unicode-math} % this also loads fontspec
  \defaultfontfeatures{Scale=MatchLowercase}
  \defaultfontfeatures[\rmfamily]{Ligatures=TeX,Scale=1}
\fi
\usepackage{lmodern}
\ifPDFTeX\else
  % xetex/luatex font selection
\fi
% Use upquote if available, for straight quotes in verbatim environments
\IfFileExists{upquote.sty}{\usepackage{upquote}}{}
\IfFileExists{microtype.sty}{% use microtype if available
  \usepackage[]{microtype}
  \UseMicrotypeSet[protrusion]{basicmath} % disable protrusion for tt fonts
}{}
\makeatletter
\@ifundefined{KOMAClassName}{% if non-KOMA class
  \IfFileExists{parskip.sty}{%
    \usepackage{parskip}
  }{% else
    \setlength{\parindent}{0pt}
    \setlength{\parskip}{6pt plus 2pt minus 1pt}}
}{% if KOMA class
  \KOMAoptions{parskip=half}}
\makeatother
\usepackage{xcolor}
\usepackage[margin=1in]{geometry}
\usepackage{longtable,booktabs,array}
\usepackage{calc} % for calculating minipage widths
% Correct order of tables after \paragraph or \subparagraph
\usepackage{etoolbox}
\makeatletter
\patchcmd\longtable{\par}{\if@noskipsec\mbox{}\fi\par}{}{}
\makeatother
% Allow footnotes in longtable head/foot
\IfFileExists{footnotehyper.sty}{\usepackage{footnotehyper}}{\usepackage{footnote}}
\makesavenoteenv{longtable}
\usepackage{graphicx}
\makeatletter
\def\maxwidth{\ifdim\Gin@nat@width>\linewidth\linewidth\else\Gin@nat@width\fi}
\def\maxheight{\ifdim\Gin@nat@height>\textheight\textheight\else\Gin@nat@height\fi}
\makeatother
% Scale images if necessary, so that they will not overflow the page
% margins by default, and it is still possible to overwrite the defaults
% using explicit options in \includegraphics[width, height, ...]{}
\setkeys{Gin}{width=\maxwidth,height=\maxheight,keepaspectratio}
% Set default figure placement to htbp
\makeatletter
\def\fps@figure{htbp}
\makeatother
\setlength{\emergencystretch}{3em} % prevent overfull lines
\providecommand{\tightlist}{%
  \setlength{\itemsep}{0pt}\setlength{\parskip}{0pt}}
\setcounter{secnumdepth}{-\maxdimen} % remove section numbering
\ifLuaTeX
  \usepackage{selnolig}  % disable illegal ligatures
\fi
\usepackage{bookmark}
\IfFileExists{xurl.sty}{\usepackage{xurl}}{} % add URL line breaks if available
\urlstyle{same}
\hypersetup{
  hidelinks,
  pdfcreator={LaTeX via pandoc}}

\author{}
\date{\vspace{-2.5em}}

\begin{document}

\section{Metadatos ICFES}\label{metadatos-icfes}

icfes: competencia: - interpretacion\_representacion nivel\_dificultad:
2 contenido: categoria: estadistica tipo: generico contexto: laboral
eje\_axial: eje4 componente: aleatorio

\section{Question}\label{question}

Una empresa farmacéutica calculó el porcentaje de efectividad de un
tratamiento para una enfermedad. Para ello, hizo cinco estudios y, de
cada estudio, se conoció el porcentaje de efectividad que tuvo el
tratamiento. De los cinco porcentajes, solo hay una moda que es 93\% y
la mediana es 98\%.

¿Es CORRECTO afirmar que la efectividad mínima mostrada en los estudios
fue de 93\%?

\subsection{Answerlist}\label{answerlist}

\begin{itemize}
\tightlist
\item
  Sí, porque la mediana debe tener dos valores mayores y dos menores, y
  la moda garantiza que los dos menores son 93\%.
\item
  No, porque la mediana es resultado de un estudio y la moda el de otros
  dos, luego alguno de los restantes podría ser menor que 93\%.
\item
  No, porque la mediana podría ser esa, por dos estudios muy cercanos a
  100\%, dos iguales y uno menor que 93\%.
\item
  Sí, porque la mediana es el resultado de uno de los estudios y la moda
  el de todos los demás, entonces el menor es 93\%.
\end{itemize}

\section{Solution}\label{solution}

La afirmación es \textbf{CORRECTA}. Analicemos por qué:

\begin{enumerate}
\def\labelenumi{\arabic{enumi}.}
\item
  \textbf{Datos y Orden:} Tenemos 5 porcentajes de efectividad. Si los
  ordenamos de menor a mayor, los llamamos \(x_1\), \(x_2\), \(x_3\),
  \(x_4\), \(x_5\).
\item
  \textbf{Mediana:} Al ser 5 datos (un número impar), la mediana es el
  valor central (el 3º). Se nos dice que la mediana es 98\%. Por lo
  tanto, \(x_3 = 98%
  \). La lista ordenada es: \(x_1\), \(x_2\), \(98%
  \), \(x_4\), \(x_5\). Esto implica que \(x_1\) es menor o igual que
  \(x_2\), que a su vez es menor o igual que \(98%
  \), y que \(98%
  \) es menor o igual que \(x_4\), que a su vez es menor o igual que
  \(x_5\).
\item
  \textbf{Moda:} La moda es el valor que más se repite. Se nos dice que
  la moda es 93\% y es \emph{única}.

  \begin{itemize}
  \tightlist
  \item
    Para que 93\% sea la moda en 5 datos, debe aparecer al menos dos
    veces.
  \item
    ¿Podría aparecer 3 o más veces? No, porque si apareciera 3 veces, al
    ordenar los datos, el valor central (la mediana, \(x_3\)) sería
    93\%, lo cual contradice que la mediana es 98\% (ya que 93\% es
    diferente de 98\%).
  \item
    Por lo tanto, 93\% aparece exactamente \emph{dos} veces.
  \item
    Al ser la \emph{única} moda, ningún otro valor puede repetirse dos o
    más veces. Esto significa que \(x_4\) es diferente de 98\%, \(x_5\)
    es diferente de 98\%, y \(x_4\) es diferente de \(x_5\).
  \end{itemize}
\item
  \textbf{Combinando Mediana y Moda:} Sabemos que la lista ordenada es
  \(x_1\), \(x_2\), \(98%
  \), \(x_4\), \(x_5\) y que dos de estos valores son 93\%.

  \begin{itemize}
  \tightlist
  \item
    Dado que 93\% (la moda) es menor que 98\% (la mediana), los dos
    valores de 93\% deben estar necesariamente en las posiciones por
    debajo de la mediana. Es decir, deben ser \(x_1\) y \(x_2\).
  \item
    La única estructura posible para la lista ordenada es: \(93%
    \), \(93%
    \), \(98%
    \), \(x_4\), \(x_5\), donde además se cumple \(98% < x_4 < x_5
    \).
  \end{itemize}
\item
  \textbf{Conclusión sobre el Mínimo:} En la lista ordenada (\(93%
  \), \(93%
  \), \(98%
  \), \(x_4\), \(x_5\)), el valor mínimo es \(x_1\), que hemos
  determinado que debe ser 93\%.
\end{enumerate}

\textbf{Tabla de datos ordenados:}

\begin{longtable}[]{@{}ccc@{}}
\toprule\noalign{}
Posición & Valor & Descripción \\
\midrule\noalign{}
\endhead
\bottomrule\noalign{}
\endlastfoot
\(x_1\) & 93\% & Moda (valor mínimo) \\
\(x_2\) & 93\% & Moda (repetido) \\
\(x_3\) & 98\% & Mediana (valor central) \\
\(x_4\) & \textgreater{} 98\% & Mayor que la mediana \\
\(x_5\) & \textgreater{} 98\% & Mayor que la mediana \\
\end{longtable}

Por lo tanto, es correcto afirmar que la efectividad mínima mostrada fue
93\%.

\section{Meta-information}\label{meta-information}

exname: Efectividad tratamiento farmaceutico ICFES extype: schoice
exsolution: 1000 exshuffle: TRUE exnumeric: FALSE extol: 0.01 expoints:
1

\end{document}
